\section{Tensor Product of \( R \)-modules}
For simplicity, consider \( R \) as commutative ring. If \( M,N,P \) are \( R \)-modules, a map \(
\vp: M\times N \to P \) is \underline{bilinear} if \( \vp(u,-), \vp(-,v) \) are \( R \)-linear maps
\( \forall \; u \in M, v\in N \).

\begin{definition}
  A \underline{tensor product} \( M \otimes_R N \) is an \( R \)-module with a bilinear map \( \vp:
  M \times N \to M \otimes_R N\), written as \( \vp(u,v) = u\otimes v \), which is universal: \(
  \forall  \) bilinear map \( \psi : M\times N \to P \), \( \exists ! \) linear map
  \[
  f: M \otimes_R N \to P \quad s.t. \; f\circ \vp = \psi
  \] 
\end{definition}

\begin{remark}
  As usual with universal properties, this characters \( M \otimes_R N \) up to a canonical
  isomorphpisms that is commutative with \( \vp \).
\end{remark}

\begin{remark}
  The proof of existence of \( M \otimes_R N \) follows as in the case \( R = \CC \): Let 
  \[
    M \otimes_R N \coloneqq  \qo{R^{(I)}}{T}
  \] 
  where \( R^{(I)} \) is the free \( R \)-module with basis \( \{e_{(u,v)} \; | \; (u,v) \in I\} \)
  and \( I = \{(u,v) \; | \; u\in M, v\in N\} \) and \( T \) is the submodule generated by all
  elements of the form 
  \[
    e_{(a_1u_1+a_2u_2,v)} - a_1 e_{(u_1,v)} - a_2e_{(u_2,v)} \quad \forall \; u_1,u_2\in M, v\in N,
    a_1,a_2 \in R 
  \] 
  and similarly w.r.t. the second variable. And define 
  \[
    \vp(u,v) \coloneqq  \overline{e_{(u,v)}}
  \] 
\end{remark}

\begin{remark}
  It follows from this description, that every element of \( M \otimes_R N \) can be written as 
  \[
    \sum_{i=1}^r a_i \overline{e_{(u_i,v_i)}} = \sum_{i=1}^r x_i \otimes y_i 
  \] 
  where \( x_i = a_i u_i, y_i = v_i \), and it is very far from unique.
\end{remark}

\begin{remark}[\textbf{General Properties}]
  
\end{remark}
